\documentclass[10pt,a4paper]{amsart}
\usepackage[margin=19mm]{geometry}
\usepackage{amsmath,amssymb,mathtools}
\usepackage[scr=rsfs,cal=euler]{mathalfa}
\usepackage{xfrac}
\usepackage[unicode,hidelinks,bookmarks=false]{hyperref}
\newcommand{\ie}{i.e.\ }
\newcommand{\cf}{cf.\ }
\newcommand{\resp}{respectively\ }
\newcommand{\Subsection}{\subsection}
\providecommand{\nobreakdash}{\nobreak\hbox{-}}
\setlength{\emergencystretch}{2em}
\multlinegap0pt
\usepackage{amstext,mathrsfs,eurosym,dsfont}
\usepackage{enumitem}
\usepackage{tikz}
\usepackage{booktabs}
\newtheorem{theorem}{Theorem}[section]
\newtheorem{lemma}[theorem]{Lemma}
\newtheorem{proposition}[theorem]{Proposition}
\newtheorem{defin}[theorem]{Definition}
\newtheorem{notation}[theorem]{Notation}
\newtheorem{assumption}[theorem]{Assumption}
\newtheorem{condition}[theorem]{Condition}
\newtheorem{remark}[theorem]{Remark}
\newtheorem{properties}[theorem]{Properties}
\newtheorem{example}[theorem]{Example}
\newtheorem{claim}[theorem]{Claim}
\numberwithin{equation}{section}
\allowdisplaybreaks
\newcommand{\charfunc}{\mathds{1}}
\newcommand{\w}{\varpi}
\newcommand{\wmf}{\varpi^{\text{mf}}}
\newcommand{\E}{\mathbb{E}}
\newcommand{\F}{\mathbb{F}}
\newcommand{\G}{\mathbb{G}}
\newcommand{\R}{\mathbb{R}}
\newcommand{\N}{\mathbb{N}}
\newcommand{\ItF}{\mathcal{F}}
\newcommand{\prob}{\mathbb{P}}
\newcommand{\setup}{(\Omega,\mathcal{F},\mathbb{P},\mathbb{F})}
\newcommand{\bsetup}{(\overline{\Omega},\overline{\mathcal{F}},\overline{\mathbb{P}},\overline{\mathbb{F}})}
\newcommand{\cadlag}{c\`{a}dl\`{a}g }
\newcommand{\Ra}{\Rightarrow}
\renewcommand{\bar}{\overline}
\renewcommand{\hat}{\widehat}
\renewcommand{\tilde}{\widetilde}
\newcommand{\blue}{\textcolor{blue}}
\newcommand{\purple}{\textcolor{purple}}
\newcommand{\red}{\textcolor{red}}
\newcommand{\avg}{\frac{1}{N}\sum_{j = 1}^N}
\newcommand{\teal}{\textcolor{teal}}
\newcommand{\mtfX}{\mathfrak{X}}
\newcommand{\mtfZ}{\mathfrak{Z}}
\newcommand{\mtfY}{\mathfrak{Y}}
\newcommand{\mtfA}{\mathfrak{A}}
\newcommand{\mtfb}{\mathfrak{b}}
\newcommand{\mtfa}{\mathfrak{a}}
\newcommand{\mtfs}{\mathfrak{s}}
\newcommand{\mtfP}{\mathfrak{Q}}
\newcommand{\mtfp}{\mathfrak{q}}
\newcommand{\X}{\mathcal{X}}
\newcommand{\Y}{\mathcal{Y}}
\newcommand{\A}{\mathcal{A}}
\newcommand{\B}{\mathcal{B}}
\begin{document}
\title[Mean-field model for pollution abatement via cap and trade m]{Mean-field model for pollution abatement via cap and trade mechanism}
\author{Ofelia Bonesini; Giacomo Lanaro}
\date{}
\maketitle
\noindent\emph{Source and licence.} Mean-field model for pollution abatement via cap and trade mechanism. Version arXiv:2607.22638v1 (CC BY 4.0). This material is distributed under the Creative Commons Attribution 4.0 International licence, http://creativecommons.org/licenses/by/4.0/, as recorded in the arXiv OAI licence metadata for this exact version and shown on the abstract page https://arxiv.org/abs/2607.22638v1. Full rights notice: licenses/arxiv-2607.22638-CC-BY-4.0.txt.\par\smallskip
\noindent\textit{Editorial scope.} Counted unit: the original \texttt{lemma} environment \texttt{lemma: filtering error independent of control} at source lines 748--753 together with its complete proof at lines 754--792; the delivered document keeps source lines 181--912 so that every referenced label and every citation resolves inside it. Native editable source is the paper's own concatenated TeX; the AMS wrapper is editorial formatting only. No publisher class, upstream script or unrelated artwork is redistributed.
\par\smallskip\noindent\textit{Reference note.} This article ships no compiled bibliography (biblatex); the entries delivered here are composed from the author's own BibTeX (.bib) fields, with no field invented and no entry dropped.
\begin{equation}
\label{eqn: empirical mean notation}
\mathfrak{m}^{(N)}(\underline{A}) := \frac{1}{N} \sum_{j = 1}^N A^j.
\end{equation}
We suppose that firm $i$ faces a linear inverse demand function defined by
\begin{equation}
\label{eqn: price good economy n}
p(K^i_t,\mathfrak{m}^{(N)}(\underline{K}_t)):= a- b(1-\gamma) AK^i_t - b\gamma A\mathfrak{m}^{(N)}(\underline{K}_t),\quad t\in[0,T],
\end{equation}
where $\underline{K}_t:= (K^1_t,\dots, K^N_t)$, $a,b>0$ and $\gamma \in[0,1]$ describes the degree of production substitution, while $A$ represents the technological level so that $AK^i_t$ is the production function of the firm $i$.

Finally, we suppose that the trading activities of firm $i$, necessary to provide the best amount of permits needed at any time is linear quadratic in the number of permits $\alpha^{x,i}_t$ traded at every time $t\in[0,T]$. The trading cost is supposed to be: 
\begin{equation}
\label{eqn: quadratic cost functional}
C^x(\alpha^{x,i}_t,\w_t) := \w_t \alpha^{x,i}_t + c_{2,x} (\alpha^{x,i}_t)^2,\quad t\in[0,T],
\end{equation}
where $c_{2,x}>0$ represents the transaction costs for trading activities and $\w_t$ is the market price of the permits at time $t\in[0,T]$. We make now the following assumption
\begin{assumption}
\label{ass: price taker companies}
In analogy to \cite{del2024mean}, we assume that no company is large enough to influence the price process $\w$. In other words, all companies are assumed to be \emph{price takers}, and $\w$ is therefore regarded as an exogenous stochastic process from the perspective of the firms, possibly correlated with all sources of randomness in the market.
\end{assumption}
To be more rigorous, we introduce the model setup. It is given by a filtered probability probability space $(\Omega,\ItF,\prob,\F)$, on which the following processes are defined:
\begin{itemize}
\item A one-dimensional Brownian motion $W^0$, playing the role of the common noise;
\item A family of $N$ idiosyncratic noises describing the randomness in the production dynamics $\underline{W}^{N;K} := (W^{K,1},\dots ,W^{K,N})$.
\item A family of $N$ idiosyncratic noises describing the randomness in the emission dynamics $\underline{W}^{N;E} := (W^{E,1},\dots ,W^{E,N})$.
\end{itemize}
We suppose that $(W^0,\underline{W}^{N;K},\underline{W}^{N;E})$ is a $(2N+1)$-dimensional Brownian motion.

For the moment, the price $\w$ is a stochastic process defined on $(\Omega,\ItF,\prob,\F)$. It is natural to suppose that each company $i$ has access to the information given by $W^0$, $\w$, and $W^i:=(W^{K,i},W^{E,i})^\top$. As a consequence, adopting the notation
\begin{align*}
v^i_t 			&:= (\alpha^{f,i}_t,\alpha^{g,i}_t,\alpha^{e,i}_t,\alpha^{x,i}_t),\quad t\in[0,T],\\
\mathcal{X}^i_t &:= (K^i_t, X^i_t), \quad t\in[0,T],
\end{align*}
for respectively the control and the state variable of firm $i = 1,\dots,N$, every firm has to minimise the following target function
\begin{align}
\label{eqn: target companies N}
&J(v^i;\mathcal{X}^i,\mathfrak{m}^{(N)}(\underline{K}), \w) \\
\nonumber& \quad:= \E\Bigg[\int_0^T \bigg[-p(K^i_t,\mathfrak{m}^{(N)}(\underline{K}_t)) AK^i_t + \alpha^{x,i}_t \w_t + c_{2,x}(\alpha^{x,i}_t)^2 + \sum_{p = f,g,e}C^p(\alpha^{p,i}_t)\bigg] dt + \lambda (X^i_T)^2\Bigg],
\end{align}
where $-p(K^i_t,\mathfrak{m}^{(N)}(\underline{K}_t)) AK^i_t$ represents the revenues coming from capital production, and $\lambda(X^i_T)^2$ is a final penalisation term present to avoid accumulation of permits in the bank accounts of the firms. In addition, we suppose that the class of admissible controls for the company $i$-th is
\begin{equation*} 
v^i \in \mathbb{A}^i := \mathbb{H}^2( \F^{W^0, W^{K,i}, W^{E,i} ,\w};\R^4),\quad i = 1,\dots, N.
\end{equation*}

Finally, the state variable of firm $i$ is fully determined by the following two dimensional SDE: 
\begin{equation*}
\begin{cases}
dK^i_t &= (a^f \alpha^{f,i}_t + a^g \alpha^{g,i}_t - \delta K^i_t) dt + \sigma^K K^i_t dW^{K,i}_t, \quad K^i_0=K^0.\\
dX^i_t &= (\alpha^{x,i}_t - a^e \alpha^{f,i}_t + \alpha^{e,i}_t) dt -\sigma^E(\rho dW^{E,i}_t + \sqrt{1- \rho^2} dW^0_t), \quad X^i_0=X^0.\\
\end{cases}
\end{equation*}
Adopting the notation of \eqref{eqn: empirical mean notation}, for every $\mathcal{X}^i\in \R^2,\  v^i\in \R^4$, for all $i = 1,\dots,N$, for all $\w\in\R$ we define:
\begin{align*}
f(\mathcal{X}^i,v^i;\mathfrak{m}^{(N)}(\mathcal{X}),\w) &:= -p(K,\mathfrak{m}^{(N)}(\underline{K})) AK^i +\sum_{p \in\{ f,g,e,x\}} C^p(\alpha^{p,i}) \\
&\qquad-\Big[ a-b(1-\gamma)AK^i - b\gamma A\mathfrak{m}^{(N)}(\underline{K})\Big]AK^i\\
&\qquad+  \sum_{p \in \{f,g,e\} } \Big[c_{1,p}\alpha^{p,i} + c_{2,p} (\alpha^{p,i})^2\Big]+ \w\alpha^{x,i} + c_{2,x} (\alpha^{x,i})^2\\
 &= \langle Q_1 \mathcal{X}^i ,\mathcal{X}^i\rangle+ 2\langle {Q}_2 \mathfrak{m}^{(N)}(\underline{\mathcal{X}}), \mathcal{X}^i\rangle + \langle Q_3 v^i, v^i\rangle + 2\langle q_1, \mathcal{X}\rangle + 2 \langle q_2(\w),v^i\rangle.\\
g(\mathcal{X}^i)&:= \langle Q_4\mathcal{X}^i, \mathcal{X}^i\rangle.
\end{align*}
where $Q_1,Q_2,Q_4\in M_{2,2}(\mathbb{R})$ are positive semi-definite, $Q_3\in M_{4,4}(\mathbb{R})$ is positive definite, $q_1\in\mathbb{R}^2$ and $q_2(\w)\in\mathbb{R}^4$, the latest depending on $\w$:
\begin{align*}
Q_1&:= 
\begin{pmatrix} 
b(1-\gamma) A^2 & 0 \\
0&0
\end{pmatrix},
\quad {Q}_2 :=
\begin{pmatrix} 
\frac{1}{2}b\gamma A^2 &0\\
0&0
\end{pmatrix},
\  Q_3:= 
\begin{pmatrix} 
c_{2,f} 	&0			&0			&0\\
0		&c_{2,g} 		&0			&0\\
0		&0			&c_{2,e}		&0\\
0		&0			&0			&c_{2,x}
\end{pmatrix},
\\
Q_4&:=
\begin{pmatrix}
 0&0\\
 0&\lambda
 \end{pmatrix} ,\ 
q_1 := 
\begin{pmatrix} 
-\frac{1}{2}aA \\
0
\end{pmatrix}
\  q_2(\w) := 
\begin{pmatrix}
\frac{1}{2}c_{1,f}\\
\frac{1}{2}c_{1,g}\\
\frac{1}{2}c_{1,e}\\
\frac{1}{2}\w
\end{pmatrix}.
\end{align*}
On the other hand, the state variable of the firm $i$ can be rewritten as in \cite{del2024mean}:
\begin{equation}
\label{eqn: state variable N}
\begin{cases}
d\mathcal{X}^i_t = \big( A_1 \mathcal{X}^i_t + A_2 v^i_t\big) dt + \sum_{j \in\{ K,E\}} (A^j_3\mathcal{X}^i_t+ A^j_4)dW^{j,i}_t + A_5dW^0_t,\\
\mathcal{X}^i_0 = (K_0,X_0),
\end{cases}\quad t\in[0,T],
\end{equation}
where for the company $i$, $W^i := (W^{K,i}, W^{E,i})$ and
\begin{align*}
A_1 &:= 
\begin{pmatrix}
-\delta	&0	\\
0		&0	
\end{pmatrix},
\ 
A_2 = 
\begin{pmatrix} 
a^f	&a^g		&0	&0\\
-a^e	&0		&1	&1
\end{pmatrix},
\ A^K_3=
\begin{pmatrix} 
\sigma^K	&0	\\
0		&0
\end{pmatrix},\\
\ A^K_4 &= \mathds{O},
\ A^E_3 = \mathds{O}, \  A^E_4 = \begin{pmatrix}
0\\
-\sigma^E\rho
\end{pmatrix},\
A_5 = \begin{pmatrix}
0\\
-\sigma^E\sqrt{1-\rho^2}
\end{pmatrix}.
\end{align*}
In analogy to \cite{del2024mean}, we denote the volatility function by
\begin{align*}
    \sigma(\mathcal{X}) &:= 
    \begin{pmatrix}
        A^K_3\mathcal{X} & A^E_4  
    \end{pmatrix}
    =
    \begin{pmatrix} 
    \sigma^K K	&0	\\
    0		& -\sigma^E \rho
    \end{pmatrix}\in M_{2\times 2}(\R) ,\\  
    \sigma^0 &:= A_5 \in M_{2\times 1}(\R).
\end{align*}
By Lemma \cite[Lemma 1.56]{carmona2018probabilistic2} and \cite[equation (6.82)]{carmona2018probabilistic2}, for all processes $\bar{\xi}, \w \in \mathbb{H}^2([0,T],\R)$, there exists a unique control $\widehat{v}^i$ such that $J(\widehat{v}^i;\mathcal{X}^i,\bar{\xi},\w)= \inf_{v \in \mathbb{H}^2(\mathbb{F}^i;\mathbb{R}^4)}J(v;\mathcal{X}^i,\bar{\xi},\w)$. 
Indeed, for every
\begin{equation*}
    \mathcal{Y}: = \begin{pmatrix}V\\ Y\end{pmatrix}\in\R^2,\quad Z:= \begin{pmatrix} Z^{1,1}&Z^{1,2}\\Z^{2,1}&Z^{2,2}\end{pmatrix}\equiv \begin{pmatrix}
        Z^K, Z^E
    \end{pmatrix}\in M_{2\times2}(\R) ,\quad Z^0 := \begin{pmatrix}Z^{0,1}\\ Z^{0,2}\end{pmatrix} \in \R^2,
\end{equation*}
the full Hamiltonian satisfies the following structure: 
\begin{align*}
&H(v,\mathcal{X}, \mathcal{Y},Z^0,Z;\bar{\xi},\w)= \\
&\qquad =  \langle A_1 \mathcal{X} + A_2v,\mathcal{Y}\rangle + f(\mathcal{X}, v; \bar{\xi}, \w)	+ \text{trace}\bigg\{\bigg( \sigma(\mathcal{X}),A_5\bigg)\begin{pmatrix}
    Z^\top\\
    (Z^0)^\top\end{pmatrix}\bigg\}\\
&\qquad= \langle A_1 \mathcal{X} + A_2v,\mathcal{Y}\rangle + f(\mathcal{X}, v; \bar{\xi}, \w)	+ \text{trace}\bigg( \sigma( \mathcal{X} ) Z^\top +  A_5(Z^0)^\top\bigg) 	\\
&\qquad= (a^f\alpha^{f} + a^g\alpha^{g}- \delta K)V + (\alpha^x -a^e\alpha^f + \alpha^e) Y + \langle Q_1 \mathcal{X} ,\mathcal{X}\rangle+ 2\langle Q_2\overline{\xi}, \mathcal{X}\rangle + \langle Q_3 v, v\rangle  \\
&\qquad\qquad+2\langle q_1, \mathcal{X}\rangle + 2 \langle q_2(\w),v\rangle + \text{trace}\bigg\{ 
\begin{pmatrix}
    \sigma^K K Z^{1,1} & \sigma^K K Z^{2,1}\\
    -\sigma^E \rho Z^{1,2} & -\sigma^E \rho Z^{2,2} 
\end{pmatrix} \\
&\qquad\qquad+ 
\begin{pmatrix}
    0&0\\
    -\sigma^E\sqrt{1-\rho^2}Z^{0,1}&-\sigma^E\sqrt{1-\rho^2}Z^{0,2}
\end{pmatrix}\bigg\}\\ 
&\qquad = (a^f\alpha^{f} + a^g\alpha^{g}- \delta K)V + (\alpha^x -a^e\alpha^f + \alpha^e) Y + \langle Q_1 \mathcal{X} ,\mathcal{X}\rangle+ 2\langle Q_2\overline{\xi}, \mathcal{X}\rangle + \langle Q_3 v, v\rangle  \\
&\qquad\qquad+2\langle q_1, \mathcal{X}\rangle + 2 \langle q_2(\w),v\rangle + \sigma^KKZ^{1,1} - \sigma^E [ \rho Z^{2,2} + \sqrt{1-\rho^2} Z^{0,2}].
\end{align*}
As a consequence, the first order condition becomes
\begin{equation}
\label{eqn: optimal control}
2Q_3 \widehat{v} + 2q_2(\w)+ A_2^\top \mathcal{Y}= 0,\Longrightarrow \widehat{v}:= -Q_3^{-1} \bigg(q_2(\w)+\frac{1}{2}A_2^\top \mathcal{Y}\bigg), 
\end{equation}
which is explicitly given by:
\begin{equation}
\label{eqn: FOC} 
\begin{cases}
\widehat{\alpha}^f &:= -\frac{1}{2c_{2,f}}(a^fV -a^eY + c_{1,f}),\\
\widehat{\alpha}^g &:=-\frac{1}{2c_{2,g}}(a^gV + c_{1,g}),\\
\widehat{\alpha}^e&:= -\frac{1}{2c_{2,e}}(c_{1,e} + Y),\\
\widehat{\alpha}^x&:= -\frac{1}{2c_{2,x}}(\w + Y).
\end{cases}
\end{equation}

In particular, the optimal control is $\widehat{v}_t=-Q_3^{-1}(q_2(\w_t)+\frac{1}{2}A_2^\top \mathcal{Y}_t)$, where $\mathcal{Y}_t$ solves the backward stochastic differential equation:


\begin{equation} 
\label{eqn: adjoint process N}
\begin{cases}
d\mathcal{Y}^i_t = -\big( A_1^\top \mathcal{Y}^i_t + (A^K_3)^\top Z^{i,K}_t + 2Q_1\mathcal{X}^i_t + 2Q_2 \mathfrak{m}^{(N)}(\underline{\mathcal{X}}_t) + 2q_1 \big)dt + Z^i_t dW^i_t + Z^{0,i}_t dW^0_t ,\\
\mathcal{Y}^i_T = Q_4\mathcal{X}^i_T.
\end{cases}
\end{equation} 

\subsection{The regulator problem}
\label{subsect_model_N:regulator}
As discussed in the introduction, the model setup considered in this paper is analogous to the one analysed in \cite[Section 3]{del2024mean}. In contrast to the cited reference, our focus is on the optimal regulatory strategy aimed at balancing carbon-emissions reduction and economic sustainability. 

The ETS mechanism is based on an auction system in which the regulator issues emission permits. The market price $\varpi$ is then determined by the interaction between firms’ demand for permits and the number of permits issued by the regulator. Consequently, by deciding over time how many \emph{new} permits to release into the market, the regulator can influence the price process and, in turn, the firms’ optimal decisions in the maximisation of their objective functionals.

%%%%%%%%%%%%%%
% AUCTION SYSTEM
%%%%%%%%%%%%%%
\subsubsection{The auction system}
We describe how the auction system works. We consider the equilibrium price, determined by the balance between the demand (number of permits requested by the companies) and the supply (number of permits sold by the companies as well as issued by the regulator in the market). This condition, called \emph{market clearing} is defined by
\begin{equation}
\label{eqn: market clearing N market}
\sum_{j = 1}^{N} \alpha^{x,j}_t + \beta^{(N)}_t = 0,\quad dt\otimes \mathbb{P},
\end{equation} 
where the control $\beta^{(N)}:= (\beta^{(N)}_t)_{t\in[0,T]}$  represents the (total) number of permits issued at time $t$ by the regulator.

\begin{remark}
\label{rmk:beta-sign}
The number of permits issued by the regulator should be of the same order of magnitude of the number of companies in the market. This is convenient to guarantee that the effect of the regulator's policies does not become negligible when the number of companies is large (for details see \cite[Remark 3.11]{fujii2022equilibrium}). As a consequence, we assume that the number of permits, in the market populated by $N$ companies, is of the form $\beta^{(N)}_t \equiv N {\beta}_t$, where $\beta_t$ denotes the \emph{normalised number of permits} issued by the regulator. 
In addition, from an economic point of view it is natural to suppose that the regulator cannot withdraw permits from the market, but it can only issue them. As a consequence, $\beta_t\leq0$ for every $t\in[0,T]$, i.e. $A= [-a,0]$ (eventually we can assume that $A = (-\infty, 0]$). 
\end{remark}

The companies are assumed to be price takers. Therefore, when solving their individual optimal control problem, they treat the price process $\w$ as exogenous. Hence, the stochastic maximum principle derived above applies and the optimal control $\hat{\alpha}^{x,j}_t$ of every company, $j \in \{1, \dots, N\}$, can be described by the last component of the first order condition on the Hamiltonian: $\widehat{\alpha}^{x,j}_t= -\frac{1}{2c_{2,x}}(\w_t + Y^j_t)$. Substituting this control in \eqref{eqn: market clearing N market}, the equilibrium price solves
\begin{equation*} 
    \sum_{j = 1}^{N} -\frac{1}{2c_{2,x}}(\w_t + Y^j_t) + \beta^{(N)}_t = 0.
\end{equation*}
As discussed in \cite{follmer1974random}, when the number of companies is sufficiently large, it is convenient to consider the normalised version of \eqref{eqn: market clearing N market}, obtained by dividing \eqref{eqn: market clearing N market} by $N$. We then have 
\begin{equation} 
\label{eqn: eq price process N}
    \w^N_t 
    \equiv -\frac{1}{N}\sum_{j = 1}^{N}Y^j_t + 2c_{2,x}\frac{1}{N}\beta^{(N)}_t 
    = -\mathfrak{m}^{(N)}(\underline{Y}_t) + 2c_{2,x}\beta_t. 
\end{equation}
This definition is in accordance with Remark \ref{rmk: controls reg N} below. Indeed, the price $\w^N$ depends on $\beta^{(N)}$ but also on 
$\mathfrak{m}^{(N)}(\underline{Y}_t)$. 
This suggests that assuming that $\beta^{(N)}$ only depends on the common shock $W^0$ would be overly restrictive. 

We adopt the notation 
\begin{equation*}
\mathfrak{m}^{(N)}(\underline{\mathcal{Y}}_t) := \mathfrak{m}^{(N)}((\underline{V}_t,\underline{Y}_t)^\top) =\frac{1}{N} \sum_{j = 1}^N (V^j_t,Y^j_t)^\top.
\end{equation*}
If the price satisfies \eqref{eqn: eq price process N}, we have an interdependent structure given by the empirical mean of the optimal controls of the companies as well as the control of the regulator. Considering the solution $\mathcal{X}$ of \eqref{eqn: state variable N}, where the control $v$ is the candidate optimal control introduced in \eqref{eqn: FOC}, we observe that the drift of $\mathcal{X}$ depends on $\w^N$. Hence, substituting $\w^N$ in the form \eqref{eqn: eq price process N} in \eqref{eqn: FOC}, we conclude that $\mathcal{X}$ depends on $\beta^{(N)}$. 


%%%%%%%%%%%%%%%%
% REGULATOR'S TARGET
%%%%%%%%%%%%%%%%

\subsubsection{The regulator's target}

In this setting, we introduce an optimal control problem for the regulator, who plays the role of a central planner. We recall that $K^j_t$ represents the capital level produced at time $t$ by company $j = 1,\dots,N$, while $E^j_t$ represents the cumulated level of emission of company $j$. Hence, the regulator is not solely interested in minimising total emissions, but rather in optimising the trade-off between environmental sustainability (i.e., reducing emissions) and industrial sustainability (i.e., supporting the economic performance of firms). 
From an economic point of view, a realistic objective regarding industrial sustainability is the maximisation of average profit of the firms, which depends on their production, energy composition, and emission abatement strategies. This approach follows the same reasoning as in \cite[equation (8)]{aid2023optimal}, where the policy maker seeks to optimise social welfare derived from firm-level performance. Accordingly, in the finite-dimensional market we define the regulator’s objective functional as
\begin{align}
\label{eqn: target-reg-N}
    J^0_{(N)}(\beta^{(N)})
    %\\
    %\nonumber&\quad
    & :=\E\Bigg[\frac{1}{N}\sum_{i = 1}^N \Bigg\{ \int_0^T a(t)\bigg[
    b(t)\beta_t^2-p(K^i_t, \mathfrak{m}^{(N)}(\underline{K}_t)) A K^i_t 
    + \hat{\alpha}^{x,i}_t \, \w^N_t + c_{2,x}(\hat{\alpha}^{x,i}_t)^2 
     \\
    \nonumber&\qquad
    + \sum_{p = f,g,e} C^p(\hat{\alpha}^{p,i}_t)
    \bigg] dt+ a(T)\lambda (X^i_T)^2 + \Bigg( \avg E^j_T - \theta\Bigg)^2 \Bigg\} \Bigg],
\end{align}

where $\w^N$ is defined in \eqref{eqn: eq price process N} and $\theta$ denotes the regulator’s target level for average economy-wide emissions after the implementation of carbon-reduction policies.

In its most straightforward formulation, the regulator’s objective combines two effects: the maximisation of firm profits and the minimisation of emissions. However, these two effects enter the functional with different growth structures. Indeed, firms’ profits depend quadratically on the production level through the revenue term, whereas cumulative emissions enter linearly. As a consequence, the incentive to sustain production tends to dominate the incentive to reduce emissions. This imbalance reflects the absence of an explicit representation of the \emph{cap} side of the cap-and-trade mechanism, i.e., the gradual limitation of the total number of permits over time. Consequently, the regulator faces a strong bias towards issuing more permits to support production, with only a relatively weaker incentive to curb emissions.
To prevent this unrealistic behaviour, we introduced a positive quadratic term $b(t)\beta_t^2$, with  $b(t) > c_{2,x}$ for all $t\in[0,T]$, which is included in the objective functional to penalize excessive new permits issuance. This term represents an implicit cost borne by the regulator when a high number of new permits is issued. Choosing $b(t)$ as an increasing function over time reinforces the interpretation that environmental policies become progressively more stringent as environmental concerns grow. In this formulation, the regulator evaluates the aggregate economic performance of the firms through their expected profits, while internalising the environmental cost generated by their emissions. The resulting policy $\beta$ determines the number of new emission permits to be issued effectively steering the trade-off between the stimulation of industrial growth and the limitation of environmental degradation.

In addition, in order to quantify the effect of the ETS-based policy of the regulator, the value $\theta$ should be defined starting from the expected amount of CO2 emitted by the economy in absence of policy, e.g. as $50\%$ of the CO2 emitted under the same economic conditions, but with no policy applied by the regulator.


\begin{remark}
\label{rmk: controls reg N}
In a market populated by finitely many companies, the strategy of the regulator should depend on any source of randomness: on the common shocks as well as the idiosyncratic shocks of every agent. However, when the number of companies is large, it is unrealistic to expect that the regulator takes her strategies in response to the noises of every company. It is actually more realistic to suppose that she can observe the common shocks affecting the economy as well as the equilibrium price process, which implicitly depends on every idiosyncratic term. Hence, the proper class of controls for the regulator should be $\mathbb{A}^{0,N} : =\mathbb{H}^2( \F^{W^0,\w^N};\R)$, where $\F^{W^0,\w^N}$ is the filtration generated by $W^0$ and $\w^N$, which is possibly correlated with the idiosyncratic noises $\underline{W}^{N;E}, \underline{W}^{N;K}$, through $\w^N$. As discussed in \cite[Remark 3.1]{fujii2022equilibrium}, this is a key issue, which leads the solution of the optimal control problem for the regulator in the finite dimensional case very difficult to solve. As we are going to discuss in the next section, a natural strategy to circumvent this problem is to pass to the mean-field formulation of such model. 
\end{remark}


Finally, we introduce the following objects to emphasise the dependence of the FBSDE on the control of the regulator:
\begin{equation*}
q_3 := 
\begin{pmatrix}
\frac{1}{2}c_{1,f}\\
\frac{1}{2}c_{1,g}\\
\frac{1}{2}c_{1,e}\\
0
\end{pmatrix},\quad 
Q_5=
\begin{pmatrix}
0	\\
0	\\
0	\\
c_{2,x}
\end{pmatrix},
\quad Q_6 :=
\begin{pmatrix}
0	&0\\
0	&0\\
0	&0\\
0	&-\frac{1}{2}\\
\end{pmatrix},
\end{equation*}
As a consequence, FBSDE \eqref{eqn: state variable N}, \eqref{eqn: adjoint process N} becomes
\begin{equation}
\label{eqn: eq FBSDE company i}
\begin{cases}
d\mathcal{X}^i_t = \big[ A_1 \mathcal{X}^i_t + A_2( -Q_3^{-1} ( q_3 + Q_5{\beta}_t + Q_6\mathfrak{m}^{(N)}(\underline{\mathcal{Y}}_t)+ \frac{1}{2}A_2^\top \mathcal{Y}^i_t))\big] dt + \sigma(\mathcal{X}^i_t) dW^i_t + A_5dW^0_t,\\
d\mathcal{Y}^i_t = -\big( A_1^\top \mathcal{Y}^i_t + (A^K_3)^\top Z^{i,K}_t + 2Q_1\mathcal{X}^i_t + 2Q_2 \mathfrak{m}^{(N)}(\underline{\mathcal{X}}_t) + 2q_1 \big)dt + Z^i_t dW^i_t + Z^{0,i}_t dW^0_t ,\\
\mathcal{X}^i_0 = (K_0,X_0)^\top,\\
\mathcal{Y}^i_T = Q_4\mathcal{X}^i_T.
\end{cases}
\end{equation}


In this formulation, the regulator’s problem is convex in the regulator's control, ensuring the existence of an optimal solution while preserving a consistent economic interpretation. 
From a mathematical standpoint, this indeed guarantees the direct application of the stochastic maximum principle as presented, for example, in \cite[Appendix A]{fujii2022equilibrium}, since the problem does satisfy the standard convexity assumptions required for optimality.

\subsection{Rewriting the regulator's objective}\label{subsect_model_N:regulator_rewritten}

Let us now analyse every term of \eqref{eqn: target-reg-N}, to rewrite more explicitly the target function of the regulator. 
In order to simplify the notation we introduce the quantity $\beta_t:=\frac{1}{N}\beta^{(N)}_t$, for all $t \in [0,T]$.
\begin{align*}
    \frac{1}{N}\sum_{j = 1}^N p(K^j_t, \mathfrak{m}^{(N)}(\underline{K}_t)) AK^j_t & = \frac{1}{N}\sum_{j = 1}^N\bigg(a- b(1-\gamma) AK^j_t - b\gamma A\mathfrak{m}^{(N)}(\underline{K}_t)\bigg) AK^j_t\\
    & = aA\mathfrak{m}^{(N)}(\underline{K}_t) - b(1-\gamma) A^2\frac{1}{N} \sum_{j = 1}^N (K^j_t)^2 - b\gamma  A^2(\mathfrak{m}^{(N)}(\underline{K}_t))^2,\\
    %%%%%
    \frac{1}{N}\sum_{j = 1}^NC^f(\hat\alpha^{f,j}_t) 
    & =\frac{1}{4 c_{2, f}}\bigg(-c_{1,f}^2 + \avg(a^f{V}^j_t-a^e{Y}^j_t)^2\bigg),\\
    %%%%%%%%%%%
    \frac{1}{N}\sum_{j = 1}^NC^e(\hat\alpha^{e,j}_t)
    &= \frac{1}{4 c_{2, e}}(-c_{1,e}^2 + \frac{1}{N}\sum_{j=1}^N(Y^j_t)^2),\\
    %%%%%%%%%%%%%%%%%
    \avg C^g(\hat\alpha^{g,j}_t) 
    &= \frac{1}{4 c_{2, g}}\bigg(-c_{1,g}^2 + (a^g)^2\frac{1}{N}\sum_{j = 1}^N (V^j_t)^2\bigg),\\
    %%%%%%%%%%%%%%%
    \avg \big( \hat \alpha^{x,j}_t \w^N_t + c_{2,x} (\hat\alpha^{x,j}_t)^2)&= \avg \Bigg( - \frac{1}{2c_{2,x}} \w^N_t(\w^N_t + Y^j_t) + \frac{1}{4c_{2,x}} (\w^N_t + Y^j_t)^2\Bigg)\\
    &= \frac{1}{4c_{2,x}}\Bigg(\avg(Y^j_t)^2  - (\w^N_t)^2\Bigg) \\
    & = \frac{1}{4c_{2,x}}\Bigg[ \avg(Y^j_t)^2 - (\mathfrak{m}^{(N)}(\underline{Y}_t))^2 + 4c_{2,x} \mathfrak{m}^{(N)}(\underline{Y}_t)\beta_t - 4(c_{2,x})^2 \beta^2_t\Bigg],\\
    %%%%%%%%%%%%%%
    \bigg(\avg E^j_T - \theta\bigg)^2 
    & = \bigg(\int_0^T \avg \hat{\alpha}^{j,i}_t dt - \avg X^j_T + A^0-\theta\bigg)^2\\
    & = \bigg (\int_0^T \beta_tdt\bigg)^2 + ( \mathfrak{m}^{(N)}(\underline{X}_T))^2 + (A^0 - \theta)^2 \\
    &\qquad + 2\bigg(\int_0^T \beta_tdt\bigg)\mathfrak{m}^{(N)}(\underline{X}_T) - 2(A^0 - \theta) \mathfrak{m}^{(N)}(\underline{X}_T)\\
    &\qquad- 2(A^0-\theta) \bigg(\int_0^T\beta_t dt\bigg).
\end{align*}

In conclusion, adopting the notation $(A^j)_{j= 1}^N$, $\mathfrak{m}^{(2,N)}(\underline{A}_t) := \frac{1}{N} \sum_{j =1}^N (A^j_t)^2$ for $t\in[0,T]$, the target function for the central planner is 
\begin{align*}
    &J^0_{(N)}(\beta^{(N)}) \\
    &= \E\Bigg[\frac{1}{N}\sum_{i = 1}^N \bigg[ \int_0^Ta(t)\bigg[b(t)\left(\frac{1}{N}\beta^{(N)}_t\right)^2 -p(K^i_t, \mathfrak{m}^{(N)}(\underline{K}_t)) AK^i_t + \hat\alpha^{x,i}_t \w^N_t + c_{2,x}(\hat\alpha^{x,i}_t)^2 \\
    &\qquad + \sum_{p = f,g,e}C^p(\hat\alpha^{p,i}_t)\bigg] dt + a(T)\lambda (X^i_T)^2 + \bigg( \avg E^j_t - \theta\bigg)^2\bigg] \Bigg]\\
    &= \E\bigg[ \int_0^Ta(t) \bigg\{-aA\mathfrak{m}^{(N)} (\underline{K}_t) + b(1-\gamma) A^2 \mathfrak{m}^{(2,N)}(\underline{K}_t) + b\gamma A^2 (\mathfrak{m}^{(N)}(\underline{K}_t))^2 \\
    &\qquad + \mathfrak{m}^{(2,N)}(\underline{Y}_t) \bigg[\frac{(a^e)^2}{4c_{2,f}} + \frac{1}{4c_{2,e}} + \frac{1}{4c_{2,x}}\bigg] +  \mathfrak{m}^{(2,N)}(\underline{V}_t) \bigg[\frac{(a^f)^2}{4c_{2,f}} + \frac{(a^g)^2}{4c_{2,g}}\bigg]\\
    &\qquad - \frac{a^fa^e}{2 c_{2,f}} \avg (V^j_tY^j_t) - \frac{1}{4c_{2,x}}(\mathfrak{m}^{(N)}(\underline{Y}_t))^2 + \mathfrak{m}^{(N)}(\underline{Y}_t)\beta_t +(b(t)- c_{2,x})\beta^2_t \bigg\}dt \\
    &\qquad+ a(T) \lambda \mathfrak{m}^{(2,N)}(\underline{X}_T) +\bigg(\int_0^T\beta_tdt\bigg)^2 + ( \mathfrak{m}^{(N)}(\underline{X}_T))^2 + (A^0 - \theta)^2 \\
    &\qquad+ 2\bigg(\int_0^T \beta_tdt\bigg)\mathfrak{m}^{(N)}(\underline{X}_T) - 2(A^0 - \theta) \mathfrak{m}^{(N)}(\underline{X}_T) - 2(A^0-\theta) \bigg(\int_0^T\beta_t dt\bigg)\bigg].
\end{align*}
We introduce the cumulated-issued credit process: 
\begin{equation}
    \label{eqn: process Bt}
    B_t := \int_0^t \beta_sds.
\end{equation}
This additional controlled process should be included in the definition of the state variable of the regulator. 


\section{A mean-field limit approximation of the market}
\label{section: mean field approach}
In this setting it seems natural to pass to the mean-field limit of the game, considering the limit in the number of companies populating the market. Indeed, the equations, that we derived for all the players of the market model, depend on the state and control of the considered player, together with the empirical mean of the solutions of \eqref{eqn: eq FBSDE company i}. 

We consider a filtered probability space $(\Omega, \ItF,\prob,\F)$ where $\F:= (\ItF_t)_{t\in[0,T]}$ is the usual augmentation of the filtration generated by a $3$-dimensional Brownian motion $(W^0,W)$, where $W := (W^E,W^K)$.



\subsection{The mean-field limit of the Nash equilibrium for the companies}
\label{subsection: mf companies game}
As in \cite{lacker2019mean}, we introduce the following random variable 
\begin{equation*}
\zeta:= (a^f,a^g,a^e,\delta,\sigma^K,\sigma^E,\rho,c_{1,f},c_{2,f},c_{1,g},c_{2,g},c_{1,e},c_{2,e},c_{1,f},c_{2,f},a,b,\gamma,A,K_0, X_0,\w,\lambda),
\end{equation*} 
that is called type vector and takes values in 
\begin{equation*}
\mathcal{Z}:= \R\times (0,\infty)^6 \times(0,1) \times (0,\infty)^{12} \times \R^2 \times (0,\infty)^2.
\end{equation*}
We call type distribution $m = \mathcal{L}(\zeta)$. 
Formally, as discussed in \cite[Section 3]{lacker2019mean}, the mean-field game is defined using the vector $\zeta$ to define a \emph{representative agent's} state variable, that solves: 
\begin{equation}
\label{eqn: mf state variable}
\begin{cases}
d\mathcal{X}_t = ( A_1 \mathcal{X}_t + A_2v_t)dt + \sigma(\mathcal{X}_t)dW_t + A_5dW^0_t,\\
\mathcal{X}_0=( K_0,X_0).
\end{cases}
\end{equation}
The target function to be minimised is
\begin{equation}
\label{eqn: mf target function}
J(v; \mathcal{X}, \bar{\mathcal{X}}, \w) := \E\bigg[\int_0^T \bigg(\langle Q_1 \mathcal{X}_t ,\mathcal{X}_t\rangle+ 2\langle Q_2 \bar{\mathcal{X}}_t, \mathcal{X}_t\rangle + \langle Q_3 v_t, v_t\rangle + 2\langle q_1, \mathcal{X}_t\rangle + 2 \langle q_2,v_t\rangle \bigg)dt + \langle Q_4\mathcal{X}_T, \mathcal{X}_T\rangle\bigg],
\end{equation}
where $\bar{\mathcal{X}}_t := (\E[K_t \mid \mathcal{F}^0_t],\E[X_t\mid \mathcal{F}^0_t])^\top$, where $(\mathcal{F}^0_t)_{t\in[0,T]}$ is the usual augmentation of the natural filtration of $W^0$. As a consequence, in analogy to \cite[Section 4]{del2024mean}, the optimal control, minimising the objective function $J(v; \mathcal{X}, \bar{K}, \w)$, is given by \eqref{eqn: optimal control}, where $\mathcal{Y}$ solves the following BSDE: 
\begin{equation} 
\label{eqn: mf adjoint process}
\begin{cases}
d\mathcal{Y}_t = -\big( A_1^\top \mathcal{Y}_t + (A^K_3)^\top Z^{K}_t + 2Q_1\mathcal{X}_t + 2Q_2 \bar{\mathcal{X}}_t + 2q_1 \big)dt + Z_t dW_t + Z^{0}_t dW^0_t ,\\
\mathcal{Y}_T = Q_4\mathcal{X}_T.
\end{cases}
\end{equation} 
In analogy to \cite[Proposition 4.2]{del2024mean}, strong existence and uniqueness of solutions of the FBSDE system
\begin{equation}
\label{eqn: mf FBSDE company}
\begin{cases}
d\mathcal{X}_t = \Big[A_1\mathcal{X}_t + A_2(-Q_3^{-1} (q_2(\w_t)+\frac{1}{2}A_2^\top \mathcal{Y}_t))\big]dt + \sigma(\mathcal{X}_t)dW_t + A_5dW^0_t,\\
\mathcal{X}_0=( K_0,X_0),\\
d\mathcal{Y}_t = -\Big[ A_1^\top \mathcal{Y}_t + (A^K_3)^\top Z^{K}_t + 2Q_1\mathcal{X}_t + 2Q_2 \bar{\mathcal{X}}_t + 2q_1 \Big]dt + Z_t dWt + Z^{0}_t dW^0_t ,\\
\mathcal{Y}_T = Q_4\mathcal{X}_T,
\end{cases}\quad t\in[0,T]
\end{equation} 
is guaranteed for every choice of the process $\w$. 

\subsection{The mean-field approximation of the optimal control problem for the regulator}\label{subsection_mfg_regulator}
We describe the mean-field approximation of optimal control problem that has to be solved by the regulator. In the mean-field setting, the regulator has to deal with a continuum of companies, as we discussed in Section \ref{subsection: mf companies game}. Defining $\overline{E}_t := \mathbb{E}[E_t|\ItF^0_t]$, the target function to be minimised by the regulator in this case is
\begin{equation*}
\mathcal{J}^0(\beta) := \mathbb{E} \bigg[ \int_0^TF( \bar{\mathcal{X}}_t,  \bar{\mathcal{Y}}_t; \beta_t) dt + G(  \bar{\mathcal{X}}_T ,\bar{\mathcal{Y}}_T; \beta_T)\bigg],
\end{equation*} 
where $\beta\in \mathbb{H}^2(\F^{W^0,\w};\R)$, $\w$ is a stochastic process defined on $(\Omega,\ItF, \prob,\F)$. With respect to the cost function defined in Section \ref{subsect_model_N:regulator}, we substitute the empirical mean of the emission of every company with the conditional expectation with respect the filtration generated by the common noise of emission of the representative company. We aim at re-writing the target function in terms of the state variable of the representative company $\mathcal{X}$ and the control of the regulator, with no explicit dependence on the emission process. 
\\
To this end, we notice that the empirical mean of emission of the companies is defined by 
\begin{equation*}
\mathfrak{m}^{(N)}(\underline{E}_t) :=\frac{1}{N} \sum_{j = 1}^N \int_0^t\alpha^{x,j}_sds - \mathfrak{m}^{(N)}(\underline{X}_t) + \mathfrak{m}^{(N)}(\underline{A}^i_0),\quad t\in[0,T].
\end{equation*}
Furthermore, as discussed in Remark \ref{rmk: controls reg N}, market clearing condition \eqref{eqn: eq price process N} is incoherent with the usual assumption made in the literature, i.e. $\w$ is adapted to $\F^0$. However, in the mean-field limit it is natural to suppose that the idiosyncratic noises cancel out, obtaining
\begin{equation}
\label{eqn: mf market clearing} 
-\beta_s := \lim_{N\to\infty} \frac{1}{N}\sum_{j = 1}^N\hat\alpha^{x,j}_s \approx \E\big[\hat{\alpha}^{x}_s |\ItF^0_s\big],\quad s\in[0,T],\quad \prob-a.s.,
\end{equation}
where $\hat\alpha^x$ is the optimal permit-trading rate of a representative company. We conclude that the number of permits emitted by the regulator should be (in the mean-field limit) independent of the idiosyncratic noises of the representative company ($W$) and adapted to $\F^0$. 

Adopting the approach developed in \cite{fujii2022mean}, we consider the mean-field limit of normalised version of the market clearing condition defined for the finite dimensional market in \eqref{eqn: eq price process N}:
\begin{equation}
\label{eqn: mf eq price}
\wmf_t :=  - \overline{Y}_t + 2c_{2,x} \beta_t,
\end{equation}
where, for every $t \in [0,T]$, $\overline{Y}_t := \E[Y_t\mid \ItF^0_t]$ and $\beta_t$ is strategy of the regulator. As a consequence, analogously to \cite[Section 6]{del2024mean}, we introduce the notation $\overline{\mathcal{Y}}_t := (\E[V_t|\ItF^0_t], \E[Y_t|\ItF^0_t])^\top$ for $t\in[0,T]$. In particular, imposing that the price process follows \eqref{eqn: mf eq price}, together with Assumption \ref{ass: price taker companies}, we obtain that the solution to the mean-field game for the representative company is described by the following FBSDE system (the adjoint process is defined as the mean-field version of \eqref{eqn: adjoint process N}):
\begin{equation}
\label{eqn: eq mf FBSDE company} 
\begin{cases}
d\mathcal{X}_t = \big[ A_1 \mathcal{X}_t + A_2( -Q_3^{-1} ( q_3  + Q_5\beta_t+ Q_6\overline{\mathcal{Y}}_t+ \frac{1}{2}A_2^\top \mathcal{Y}_t))\big] dt + \sigma(\mathcal{X}_t)dW_t + A_5dW^0_t,\\
d\mathcal{Y}_t = -\big( A_1^\top \mathcal{Y}_t + (A^K_3)^\top Z^{K}_t + 2Q_1\mathcal{X}_t + 2Q_2 \bar{\mathcal{X}}_t) + 2q_1 \big)dt + Z_t dWt + Z^{0}_t dW^0_t ,\\
\mathcal{X}_0 = (K_0,X_0)^\top,\\
\mathcal{Y}_T = Q_4\mathcal{X}_T.
\end{cases}
\end{equation}
We should prove the existence and uniqueness of the solution of \eqref{eqn: eq mf FBSDE company} for every given process $\beta\in\mathbb{H}^2(\F^0;\R)$. For the existence and uniqueness of solution of FBSDEs of the Mckean-Vlasov type (i.e. whose coefficients depend on the distribution of the solution itself), we base on \cite{carmona2013mean,carmona2015forward}.

\subsection{The passage to the limit}\label{subsect_MFG_passage_to_limit}
By Yamada-Watanabe representation theorem \cite[Theorem 1.33]{carmona2018probabilistic2}, there exists a deterministic progressively measurable function $\Psi$ such that
\begin{align*}
\Psi:\mathcal{D}([0,T],\R) \times \mathcal{C}([0,T],\R^2)&\to \mathcal{C}([0,T], \R^5)\times  \mathcal{D}([0,T],\R^2) \\
(\beta, W^0, W) & \mapsto\Psi( \beta, W^0,W):= (\mathcal{X}, \mathcal{Y}, Z^0,Z),
\end{align*} 
where $(\mathcal{X},\mathcal{Y},Z^0,Z)$ solves \eqref{eqn: eq mf FBSDE company}. Reasoning analogously, there exists a progressively measurable function $\Psi^N$, such that denoted by $\mathcal{K}^{i,N}:= (\mathcal{X}^i, \mathcal{Y}^i, Z^{0,i}, Z^i)$ the solution of \eqref{eqn: eq FBSDE company i}, for $i = 1,\dots,N$, the following relation holds
\begin{align*} 
\Psi^N: \mathcal{D}([0,T],\R) \times \mathcal{C}([0,T],\R^{N+1})&\to (\mathcal{C}([0,T],\R^4))^N\\ 
 ( \beta, W^0,W^1,\dots,W^N)& \mapsto \Psi^N(\mathcal{K}^{1,N},\dots, \mathcal{K}^{N,N}).
\end{align*}
We conclude that the family $(\mathcal{K}^{1,N},\dots, \mathcal{K}^{N,N})$ is exchangeable. As a consequence, the families 
\begin{equation*}
((K^1)^2,\dots, (K^N)^2),\quad ((Y^1)^2,\dots, (Y^N)^2),\quad ((V^1)^2,\dots, (V^N)^2),\quad(V^1Y^1,\dots, V^N,Y^N)
\end{equation*}
are exchangeable too. In analogy to Section \ref{subsection: mf companies game}, we conclude that
\begin{align*} 
\lim_{N\to \infty} \mathfrak{m}^{(2,N)}(\underline{L}_t) &= \E[L_t^2 \mid \ItF^0_t],\quad t\in[0,T],\quad \prob-a.s,\quad \forall L\in \{K,X,Y,V\},\\
\lim_{N\to \infty} \avg (V^j_tY^j_t)  &= \E[V_tY_t \mid \ItF^0_t],\quad t\in[0,T],\quad \prob-a.s,
\end{align*}
where $(\mathcal{X}, \mathcal{Y}) := ((K,X), (V,Y))$ solves \eqref{eqn: eq mf FBSDE company}. In the mean-field limit the optimal control problem to be solved by the regulator becomes
\begin{align}
    \label{eqn: mf target function regulator}
    &\mathcal{J}^0(\beta) \\
    &\nonumber\qquad :=\E \Bigg[ \int_0^T\Bigg[a(t) \Bigg(  -aA\E[K_t \mid \ItF^0_t] + b(1-\gamma)  A^2\E[ K^2_t \mid \ItF^0_t] + b \gamma A^2 \Big( \E[K_t\mid \ItF^0_t] \Big)^2  \\
    &\nonumber\qquad \qquad+ \gamma_1 \E[Y^2_t\mid\ItF^0_t] + \gamma_2 \Big(\E[Y_t \mid \ItF^0_t]\Big)^2 + \gamma_3 \E[ V^2_t\mid \ItF^0_t]  + \gamma_4 \E[V_t Y_t \mid \ItF^0_t] + \gamma_5\Bigg) \\
    &\nonumber\qquad\qquad  + \beta_t \big( a(t)\E[Y_t \mid \ItF^0_t] - 2(A^0-\theta)\big) +a(t)(b(t)-c_{2,x}) \beta^2_t \Bigg] dt + a(T) \lambda \E[X^2_T \mid \ItF^0_t] + B^2_T  \\
    &\nonumber\qquad\qquad + (\E[X_T\mid \ItF^0_T])^2+ (A^0 -\theta)^2 + 2 B_T \E[X_T \mid \ItF^0_T] - 2 (A^0 - \theta) \E[X_T \mid \ItF^0_T]\Bigg]\\
    &\nonumber\qquad =\E \Bigg[ \int_0^T\Bigg[a(t) \Bigg(  -aAK_t + b(1-\gamma) A^2 K^2_t + b  \gamma A^2\Big( \E[K_t\mid \ItF^0_t] \Big)^2 + \gamma_1Y^2_t+ \gamma_2 \Big(\E[Y_t \mid \ItF^0_t]\Big)^2 \\
    &\nonumber\qquad \qquad +\gamma_3 V^2_t + \gamma_4 V_t Y_t  + \gamma_5\Bigg) + \beta_t \big( a(t)\E[Y_t\mid \ItF^0_t]  - 2(A^0 - \theta)\big) +a(t)(b(t)-c_{2,x}) \beta^2_t \Bigg] dt \\
    &\nonumber \qquad\qquad+ a(T) \lambda X^2_T+B^2_T + (\E[X_T\mid \ItF^0_T])^2 + 2B_TX_T - 2(A^0-\theta)X_T\Bigg] + (A^0-\theta)^2\\
    &\nonumber\qquad = \E\Bigg[ \int_0^T (\langle Q^{0;X}_s \mathcal{X}_s,\mathcal{X}_s\rangle+\langle Q^{0;Y}_s \mathcal{Y}_s,\mathcal{Y}_s\rangle + \langle \bar{Q}^{0;X}_s \bar{\mathcal{X}}_s, \bar{\mathcal{X}}_s\rangle+ \langle \bar{Q}^{0;Y}_s \bar{\mathcal{Y}}_s, \bar{\mathcal{Y}}_s\rangle +a(s)(b(s)- c_{2,x})\beta_s^2\\
    &\nonumber\qquad \qquad  -2(A^0-\theta)\beta_s+ 2 S^{0;Y}_s \bar{\mathcal{Y}}_s \beta_s  + 2 \langle q^{0;X}_s , \mathcal{X}_s\rangle + 2 \langle q^{0;Y}_s ,\mathcal{Y}_s\rangle + \gamma_5) ds + \langle H^0 \mathcal{X}_T , \mathcal{X}_T\rangle\\
    &\nonumber\qquad\qquad + \langle H^1\bar{\mathcal{X}}_T, \bar{\mathcal{X}}_T\rangle + 2B_T \langle h^1, \bar{\mathcal{X}}_T \rangle+ B^2_T + \langle h^0, \mathcal{X}_T\rangle\Bigg]+ (A^0 - \theta)^2,
\end{align} 
where 
\begin{alignat*}{4}
    \gamma_1 &= \frac{1}{4c_{2,x}} + \frac{(a^e)^2}{4c_{2,f}} + \frac{1}{4c_{2,e}}, \qquad&
    Q^{0;X}_s &= a(s) \begin{pmatrix} b(1-\gamma)A^2 & 0 \\ 0 & 0 \end{pmatrix}, \\[6pt]
    \gamma_2 &= -\frac{1}{4c_{2,x}}, \qquad&
    \bar{Q}^{0;X}_s &= a(s) \begin{pmatrix} b\gamma A^2 & 0 \\ 0 & 0 \end{pmatrix}, \\[6pt]
    \gamma_3 &= \frac{(a^f)^2}{4c_{2,f}} + \frac{(a^g)^2}{4c_{2,g}}, \qquad&
    Q^{0;Y}_s &= a(s) \begin{pmatrix} \gamma_3 & \frac{1}{2}\gamma_4 \\ \frac{1}{2}\gamma_4 & \gamma_1 \end{pmatrix}, \\[6pt]
    \gamma_4 &= -\frac{a^f a^e}{2c_{2,f}}, \qquad&
    \bar{Q}^{0;Y}_s &= a(s) \begin{pmatrix} 0 & 0 \\ 0 & \gamma_2 \end{pmatrix}, \\[6pt]
    \gamma_5 &= -\left[ \frac{(c_{1,f})^2}{4c_{2,f}} + \frac{(c_{1,g})^2}{4c_{2,g}} + \frac{(c_{1,e})^2}{4c_{2,e}} \right], \qquad&
    S^{0;Y}_s &= \begin{pmatrix} 0 & \frac{1}{2}a(s) \end{pmatrix}, \\[6pt]
    & &
    q^{0;X}_s &= \begin{pmatrix} -\frac{1}{2}a(s)aA \\ 0 \end{pmatrix}, \\[6pt]
    & &
    H^0 &= \begin{pmatrix} 0 & 0 \\ 0 & a(T)\lambda \end{pmatrix}, \\[6pt]
    & &
    H^1 &= \begin{pmatrix} 0 & 0 \\ 0 & 1 \end{pmatrix}, \\[6pt]
    & &
    h^0 &= \begin{pmatrix} 0 \\ -2(A^0 - \theta) \end{pmatrix}, \\[6pt]
    & &
    h^1 &= \begin{pmatrix} 0 \\ 1 \end{pmatrix}.
\end{alignat*}

The second identity comes from the tower property, Fubini's theorem together with the fact that $\beta\in \mathbb{H}^2(\F^0;\R)$. 

The regulator’s optimisation problem in the mean-field limit, given by \eqref{eqn: mf target function regulator}, is a linear–quadratic stochastic control problem of McKean–Vlasov type. In the following we omit the term $(A^0 - \theta)^2$, appearing in the cost functional, since it doesn't depend either on the control $\beta$ or the state variable $(\mathcal{X}, B, \mathcal{Y})$. The state variables satisfy a forward–backward system of stochastic differential equations driven by idiosyncratic and common sources of uncertainty, with the solution indicated by $(\mathcal{X}, B, \mathcal{Y}, Z^0, Z)$. 


In conclusion, the mean-field formulation of the optimal control problem of the regulator is given by $\inf_{\beta \in \mathbb{H}^2(\F^0;[0,\infty))}\mathcal{J}^0(\beta)$, where
\begin{align}
&\label{eqn: convex target function mf regulator}
    \mathcal{J}^0(\beta) \\
    &\nonumber\qquad = \E\Bigg[ \int_0^T (\langle Q^{0;X}_s \mathcal{X}_s,\mathcal{X}_s\rangle+\langle Q^{0;Y}_s \mathcal{Y}_s,\mathcal{Y}_s\rangle + \langle \bar{Q}^{0;X}_s \bar{\mathcal{X}}_s, \bar{\mathcal{X}}_s\rangle+ \langle \bar{Q}^{0;Y}_s \bar{\mathcal{Y}}_s, \bar{\mathcal{Y}}_s\rangle +a(s)(b(s) -  c_{2,x})( \beta_s)^2\\
    &\nonumber\qquad \qquad  -2(A^0 - \theta)\beta_s+ 2 \langle S^{0;Y}_s \bar{\mathcal{Y}}_s , \beta_s \rangle + 2 \langle q^{0;X}_s , \mathcal{X}_s\rangle + 2 \langle q^{0;Y}_s ,\mathcal{Y}_s\rangle + \gamma_5) ds + \langle H^0 \mathcal{X}_T , \mathcal{X}_T\rangle\\
    &\nonumber\qquad\qquad+ \langle h^0, \mathcal{X}_T\rangle+ \langle H^1\bar{\mathcal{X}}_T, \bar{\mathcal{X}}_T\rangle + 2B_T\langle h^1,\bar{\mathcal{X}}_T\rangle + B^2_T\Bigg],
\end{align}
subject to \eqref{eqn: eq mf FBSDE company}. 

\begin{remark}[On the mean-field formulation of the problem]
\label{rmk: defence of mfg}
One might wonder whether the model could be studied directly in a finite-dimensional setting, that is, for an economy populated by a finite number $N$ of firms. Although this is in principle possible, the analysis becomes considerably more intricate due to the impact of idiosyncratic noises on the equilibrium price and, consequently, on the agents’ strategies. 

In the $N$-player game, the equilibrium price $\w^N$ defined in \eqref{eqn: eq price process N} depends on the optimal controls of all firms and on the regulator’s policy. Since each firm’s optimal control reacts to its own idiosyncratic noise, the price $\w^N$ inherits a dependence on all firms’ individual shocks. As a result, the control of any given firm becomes correlated with the idiosyncratic noises of all other firms through a recursive structure that is extremely difficult to characterise. 

Passing to the mean-field limit resolves this difficulty. In the mean-field framework, the equilibrium price $\wmf$ becomes adapted to the common noise $\F^0$ only, as the aggregate effect of idiosyncratic noises cancels out in the limit. Consequently, each representative firm’s optimal control is adapted to its own idiosyncratic noise and to the common noise, but not to the idiosyncratic shocks of others. This greatly simplifies the structure of the equilibrium and enhances tractability.

A similar argument applies to the regulator’s problem. In the finite-dimensional case, the regulator realistically observes the common noise and the equilibrium price $\w^N$, which is itself correlated with the idiosyncratic noises of all firms. Therefore, the regulator’s control is adapted to the filtration $\mathcal{F}^{W^0, \w}$, where $\w$ depends recursively on the regulator’s own optimal policy. This feedback loop makes it extremely difficult to obtain explicit characterisations of the equilibrium. In the mean-field limit, however, the equilibrium price depends only on the common noise, so the regulator’s control becomes adapted to $\F^0$ alone, restoring analytical tractability.

Finally, it is worth noting that the mean-field formulation naturally accommodates a \emph{multi-population} extension, in which firms belong to different groups with heterogeneous technological or environmental characteristics, for instance, industries that rely more or less heavily on emissions to sustain their business models. In such a setting, one may reasonably assume that each population faces its own representative control problem within the same aggregate economy. As long as populations do not share additional sources of common information inaccessible to the others, the mean-field limit ensures that idiosyncratic noises continue to average out at the population level. Hence, the analytical simplifications afforded by the mean-field framework remain valid, and the extension to multiple populations can be obtained straightforwardly.
\end{remark}
 








\section{Analysis of MF optimal control problem subject to controlled FBSDE}\label{section_solution_MF_regulator_cpo}

The aim of this section is to study the solutions of the optimal control problem defined in Equations \eqref{eqn: eq mf FBSDE company} and \eqref{eqn: convex target function mf regulator}. This is a partial-information linear quadratic optimal control problem, where the state variable solves a controlled FBSDE system. We notice that the target function can be rewritten as 
\begin{align*}
    \mathcal{J}^0(\beta)    &= \bar{\mathcal{J}}^0(\beta) + \E\bigg[\int_0^T\bigg( \langle Q^{0,X}(\mathcal{X}_s - \bar{\mathcal{X}}_s), \mathcal{X}_s - \bar{\mathcal{X}}_s\rangle + \langle Q^{0,Y}(\mathcal{Y}_s - \bar{\mathcal{Y}}_s), \mathcal{Y}_s - \bar{\mathcal{Y}}_s\rangle \bigg) ds \\
    &\qquad+ \langle H^0 ( \mathcal{X}_T - \bar{\mathcal{X}}_T), \mathcal{X}_T - \bar{\mathcal{X}}_T\rangle \bigg],
\end{align*}
where 
\begin{align}
    &\label{eqn: convex target function mf regulator projected}
    \bar{\mathcal{J}}^0(\beta) \\
    &\nonumber\qquad = \E\Bigg[ \int_0^T \bigg(\langle (Q^{0;X}_s + \bar{Q}^{0;X}_s) \bar{\mathcal{X}}_s,\bar{\mathcal{X}}_s\rangle+\langle (Q^{0;Y}_s + \bar{Q}^{0;Y}_s) \bar{\mathcal{Y}}_s,\bar{\mathcal{Y}}_s\rangle +a(t)(b(s) -  c_{2,x})( \beta_s)^2\\
    &\nonumber\qquad \qquad-2(A^0 - \theta)\beta_s  + 2  S^{0;Y}_s \bar{\mathcal{Y}}_s  \beta_s  + 2 \langle q^{0;X}_s , \bar{\mathcal{X}}_s\rangle + 2 \langle q^{0;Y}_s ,\bar{\mathcal{Y}}_s\rangle + \gamma_5\bigg) ds + \langle H^0  \bar{\mathcal{X}}_T , \bar{\mathcal{X}}_T\rangle \\
    &\nonumber\qquad\qquad + \langle h^0, \bar{\mathcal{X}}_T\rangle+ \langle H^1\bar{\mathcal{X}}_T, \bar{\mathcal{X}}_T\rangle + 2B_T \langle h^1, \bar{\mathcal{X}}_T\rangle + B^2_T\rangle\Bigg].
\end{align}
We prove now the following result

\begin{lemma}
\label{lemma: filtering error independent of control}
\begin{equation*}\E\bigg[\int_0^T\bigg( \langle Q^{0,X}(\mathcal{X}_s - \bar{\mathcal{X}}_s), \mathcal{X}_s - \bar{\mathcal{X}}_s\rangle + \langle Q^{0,Y}(\mathcal{Y}_s - \bar{\mathcal{Y}}_s), \mathcal{Y}_s - \bar{\mathcal{Y}}_s\rangle \bigg) ds + \langle H^0  (\mathcal{X}_T - \bar{\mathcal{X}}_T), \mathcal{X}_T - \bar{\mathcal{X}}_T\rangle \bigg]
\end{equation*}
does not depend on $\beta$.
\end{lemma}

\begin{proof}
We recall that the solution $(\mathcal{X},\mathcal{Y},Z^0,Z)$ of the FBSDE \eqref{eqn: eq mf FBSDE company} is adapted to $\F$ since $\beta\in \mathbb{H}^2(\F^0;[0,\infty))$. In addition, let us notice that $\ItF^0_t = \ItF^0_s \vee \sigma\{W^0_u - W^0_s:\ u\in[0,t]\}$, for every $s\in[0,t]$, for every $t\in [0,T]$. Hence, for $s\leq t$, $\E[ \mathcal{X}_s|\ItF^0_t] = \E[ \mathcal{X}_s|\ItF^0_s] = \bar{\mathcal{X}}_s$ and the same holds for $\mathcal{Y}$, $Z^0$ and $Z$. In particular, we have
\begin{align*}
    \bar{\mathcal{X}}_t &= \mathcal{X}^0+ \E\bigg[\int_0^t \big( A_1 \mathcal{X}_s + A_2( -Q_3^{-1} ( q_3  + Q_5\beta_s+ Q_6\overline{\mathcal{Y}}_s+\frac{1}{2} A_2^\top \mathcal{Y}_s)) \big) ds \mid \ItF^0_t\bigg]+ A_5W^0_t\\
        &=  \mathcal{X}^0 + \int_0^t \big( A_1 \bar{\mathcal{X}}_s + A_2( -Q_3^{-1} ( q_3  + Q_5\beta_s+ (Q_6+\frac{1}{2}A_2^\top) \bar{\mathcal{Y}}_s)) \big) ds+ A_5W^0_t
\end{align*}
Reasoning similarly to \cite[Theorem 2.2]{wang2018introduction}, which is based on \cite[Theorem 8.1]{liptser2013statistics}, we have that
\begin{align*}
\bar{\mathcal{Y}}_t &=\mathcal{Y}_0 -\E\bigg[\int_0^t\big( A_1^\top \mathcal{Y}_s + (A^K_3)^\top Z^{K}_t + 2Q_1\mathcal{X}_s + Q_2\overline{\mathcal{X}}_s + 2q_1\big) ds - \int_0^tZ_s dW_s\\
&\qquad - \int_0^tZ^{0}_s dW^0_s\biggm| \ItF^0_t\bigg]\\
&= \mathcal{Y}_0- \int_0^t\big( A_1^\top \bar{\mathcal{Y}}_s + (A^K_3)^\top \bar{Z}^{K}_t + 2(Q_1+Q_2+q_1)\bar{\mathcal{X}}_s \big) ds + \int_0^t\bar{Z}^0_s dW^0_s.
\end{align*}
In conclusion, $(\bar{\mathcal{X}}, \bar{\mathcal{Y}}, \bar{Z}, \bar{Z}^0)$ solves the following FBSDE system 
\begin{equation}
\label{eqn: eq mf FBSDE company projected} 
\begin{cases}
    d\bar{\mathcal{X}}_t&=   \big[ A_1 \bar{\mathcal{X}}_t + A_2\big( -Q_3^{-1} ( q_3  + Q_5\beta_t+( Q_6 +\frac{1}{2} A_2^\top)\overline{\mathcal{Y}}_t)\big)\big] dt  + A_5dW^0_t,\\
    d\bar{\mathcal{Y}}_t &= - \big[A^\top_1 \bar{\mathcal{Y}}_t+ (A^K_3)^\top Z^{K}_t + 2(Q_1 + Q_2 +q_1) \bar{\mathcal{X}}_t \big] dt+ \bar{Z}^0_t dW^0_t, \\
\bar{\mathcal{X}}_0 &= (K_0,X_0)^\top,\\
\bar{\mathcal{Y}}_T &= Q_4\bar{\mathcal{X}}_T.
\end{cases}
\end{equation}

It is sufficient now to show that the filtering error,
\begin{equation*}
(\hat{\mathcal{X}} , \hat{\mathcal{Y}}, \hat{Z}) := 
(\mathcal{X} - \bar{\mathcal{X}}, \mathcal{Y} - \bar{\mathcal{Y}}, Z - \bar{Z}),
\end{equation*}
does not depend on $\beta$. We notice that $(\hat{\mathcal{X}} , \hat{\mathcal{Y}})$ solve the following FBSDE system
\begin{equation*}
    \begin{cases}
d\hat{\mathcal{X}}_t&=   \big( A_1 \hat{\mathcal{X}}_t - \frac{1}{2}A_2Q_3^{-1}  A_2^\top\hat{\mathcal{Y}}_t\big) dt  + A^K_3\hat{\mathcal{X}}_tdW_t,\\
d\hat{\mathcal{Y}}_t &= - (A^\top_1 \hat{\mathcal{Y}}_t+ (A^K_3)^\top \hat{Z}^{K}_t+ 2Q_1 \hat{\mathcal{X}}_t )dt+ \hat{Z}_t dW_t+ \hat{Z}^0_t dW^0_t, \\
\hat{\mathcal{X}}_0 &= 0,\\
\hat{\mathcal{Y}}_T &= Q_4\hat{\mathcal{X}}_T.
    \end{cases}
\end{equation*}
If uniqueness of strong solutions holds, by Yamada-Watanabe representation theorem \cite[Theorem 1.33]{carmona2018probabilistic2}, $ (\hat{\mathcal{X}}, \hat{\mathcal{Y}}, \hat{Z}^0, \hat{Z})$ is uniquely determined by $(W^0,W)$ and, in addition, does not depend on $\beta$. 
\end{proof}

In conclusion, we can reduce the regulator's problem to $\inf_{\beta \in \mathbb{H}^2(\F^0; \R)}\bar{\mathcal{J}}^0(\beta)$, introduced in \eqref{eqn: convex target function mf regulator projected} subject to $(\bar{\mathcal{X}},\bar{\mathcal{Y}})$, solution of \eqref{eqn: eq mf FBSDE company projected}. 

\subsection{Stochastic maximum principle in the regulator framework}

We start by recalling the FBSDEs system that defines the dynamics for the state: 
\begin{equation}
\label{eqn: fbsde mf repr company}
\begin{cases}
d\mathcal{X}_t 
    &= [A_{\mathcal{X}} \mathcal{X}_t + A_{\mathcal{Y}} \mathcal{Y}_t + \overline{A}_{\mathcal{Y}} \overline{{\mathcal{Y}}}_t + A_\beta \beta_t + a_\mathcal{X}]\, dt + [C_{\mathcal{X}} \mathcal{X}_t + c_K]\,dW_t^{K} +c_X\,dW_t^{X} + c_0\,dW_t^{0}, \\
    dB_t &= \beta_t dt,\\
    d{\mathcal{Y}}_t &= [B_{\mathcal{X}} \mathcal{X}_t + \overline{B}_{\mathcal{X}} \overline{{\mathcal{X}}}_t + B_{\mathcal{Y}} {\mathcal{Y}}_t + B_Z {Z}^K_t + b_\mathcal{Y}]\,dt  + {Z}^K_t\,d{W}_t^{K} + {Z}^X_t\,dW_t^{X}+ {Z}^0_t\,dW_t^{0},\\
    {\mathcal{X}}_0 &= (K_0,X_0)^\top,\\
    {\mathcal{Y}}_T &= Q_4{\mathcal{X}}_T.
\end{cases}
\end{equation}
where
\begin{alignat*}{5}
    A_{\mathcal{X}} &= A_1, &\quad
    A_{\mathcal{Y}} &= -\tfrac{1}{2}A_2 Q_3^{-1} A_2^\top, &\quad
    \overline{A}_{\mathcal{Y}} &= -A_2Q_3^{-1}Q_6, &\quad
    A_\beta &= -A_2Q_3^{-1}Q_5, &\quad
    a_{\mathcal{X}} &= -A_2 Q_3^{-1}q_3, \\[6pt]
    B_{\mathcal{X}} &= -2Q_1, &\quad
    \overline{B}_{\mathcal{X}} &= -2Q_2, &\quad
    B_{\mathcal{Y}} &= -A_1^\top, &\quad
    B_Z &= -A_3^K,&\quad
    b_{\mathcal{Y}} &= -2q_1, \\[6pt]
    C_{\mathcal{X}} &= A_3^K, &\quad
    c_K &= 0, &\quad
    c_X &= A_4^E, &\quad
    c_0 &= A_5.
\end{alignat*}
The state variable is now $(\bar{\mathcal{X}}_t, B_t)$. Nevertheless, because the process $B$ does not affect the dynamics of either $\bar{\mathcal{X}}$ or $\mathcal{Y}$, we consider $\bar{\mathcal{X}}$ and $B$ as distinct components.


The orthogonal decomposition of the states satisfies the following dynamics
\begin{align*}
    d\overline{\mathcal{X}}_t 
    &= [A_{\mathcal{X}} \overline{{\mathcal{X}}}_t + A_{\mathcal{Y}} \overline{\mathcal{Y}}_t + \overline{A}_{\mathcal{Y}} \overline{{\mathcal{Y}}}_t + A_\beta \beta_t + a_\mathcal{X}]\,dt + c_0\,dW_t^{0} \\
    &= [A_{\mathcal{X}} \overline{{\mathcal{X}}}_t + (A_{\mathcal{Y}} + \overline{A}_{\mathcal{Y}}) \overline{{\mathcal{Y}}}_t + A_\beta \beta_t + a_\mathcal{X}]\,dt + c_0\,dW_t^{0} \\
    d({\mathcal{X}} - \overline{{\mathcal{X}}})_t 
    &= [A_{\mathcal{X}}({\mathcal{X}} - \overline{{\mathcal{X}}})_t + A_{\mathcal{Y}}({\mathcal{Y}} - \overline{{\mathcal{Y}}})_t]\,dt + (C_{\mathcal{X}}({\mathcal{X}} - \overline{{\mathcal{X}}})_t + C_{\mathcal{X}} \overline{{\mathcal{X}}}_t +c_K )\,dW_t^{K} + c_X\,dW_t^{X} \\
    d\overline{\mathcal{Y}}_t 
    &= [(B_{\mathcal{X}} + \overline{B}_{\mathcal{X}}) \overline{\mathcal{X}}_t + B_{\mathcal{Y}} \overline{\mathcal{Y}}_t + B_Z \overline{Z}^K_t + b_\mathcal{Y}]\,dt + {Z}_t^{0}\,dW_t^{0} \\
    d({\mathcal{Y}} - \overline{{\mathcal{Y}}})_t 
    &= [B_{\mathcal{X}}({\mathcal{X}} - \overline{{\mathcal{X}}})_t + B_{\mathcal{Y}}({\mathcal{Y}} - \overline{{\mathcal{Y}}}) + B_Z({Z}^K_t - \overline{Z}^K_t)]\,dt + {Z}_t^{K}\,dW_t^{K} + {Z}_t^{X}\,dW^E_t
\end{align*}

Assuming the following ansatz for the form of the adjoint variable $\mathcal{Y}$
\begin{equation}\label{eq:ANSATZ_Y}
    {\mathcal{Y}}_t 
    = {P}_t ({\mathcal{X}}_t - \overline{{\mathcal{X}}}_t) + Q_t \overline{{\mathcal{X}}}_t +\Psi_t+ \varphi_t,
\end{equation}
with $P$, $Q$ and $\varphi$ deterministic function of time, such that $P_T=Q_4=Q_T$ and $\varphi_T=0$, and
\begin{align*}
    \Psi_t := \Psi_T-\int_t^T \psi_s ds + \int_t^T Z^\Psi_s dW^0_s,\quad \Psi_T = 0
\end{align*} 
with $\psi$ and $Z^\Psi$ stochastic processes adapted to $\mathbb{F}^0$, an application of the Ito formula yields
\begin{align*}
    d{\mathcal{Y}}_t 
    & = \left[ \dot{P}_t ({\mathcal{X}}_t - \overline{{\mathcal{X}}}_t) + {P}_t A_{\mathcal{X}} ({\mathcal{X}}_t - \overline{{\mathcal{X}}}_t) + {P}_t A_{\mathcal{Y}} ({\mathcal{Y}}_t - \overline{{\mathcal{Y}}}_t)\right]dt + \left[{P}_t C_{\mathcal{X}}({\mathcal{X}}_t - \overline{{\mathcal{X}}}_t) + {P}_t C_{\mathcal{X}} \overline{{\mathcal{X}}}_t +P_tc_K\right]dW_t^{K} \\
    & \quad + P_tc_X\,dW_t^{X} +\left[ \dot{Q}_t\overline{{\mathcal{X}}}_t + Q_t A_{\mathcal{X}} \overline{{\mathcal{X}}}_t + Q_t (A_{\mathcal{Y}} + \overline{A}_{\mathcal{Y}})\overline{{\mathcal{Y}}}_t + Q_t A_\beta \beta_t + Q_t a_\mathcal{X}\right]dt \\
    & \quad + (Q_t c_{0}-Z^\Psi_t)\,dW_t^{0} +\psi_t dt+ \dot{\varphi}_t\,dt.
\end{align*}
From the ansatz in \eqref{eq:ANSATZ_Y}, we derive 
\begin{equation}\label{eq:ansatz_cons}
    \overline{{\mathcal{Y}}}_t 
    = Q_t \overline{{\mathcal{X}}}_t + \Psi_t + \varphi_t
    \qquad 
    {\mathcal{Y}}_t - \overline{{\mathcal{Y}}}_t 
    = {P}_t({\mathcal{X}}_t - \overline{{\mathcal{X}}}_t)
\end{equation}
and so, replacing these expression in the previous equation, we get:

\begin{align}\label{eq:iniz_coef_dY}
    d{\mathcal{Y}}_t & = \left[ \dot{P}_t ({\mathcal{X}}_t - \overline{{\mathcal{X}}}_t) + {P}_t A_{\mathcal{X}} ({\mathcal{X}}_t - \overline{{\mathcal{X}}}_t) + {P}_t A_{\mathcal{Y}} {P}_t({\mathcal{X}}_t - \overline{{\mathcal{X}}}_t)\right]\,dt \\
    \nonumber
    & \quad + \left[{P}_t C_{\mathcal{X}}({\mathcal{X}}_t - \overline{{\mathcal{X}}}_t) + {P}_t C_{\mathcal{X}} \overline{{\mathcal{X}}}_t +P_tc_K\right]\,dW_t^{K} +P_tc_X\,dW_t^{X} \\
    \nonumber
    &\quad +\left[ \dot{Q}_t\overline{{\mathcal{X}}}_t + Q_t A_{\mathcal{X}} \overline{{\mathcal{X}}}_t + Q_t (A_{\mathcal{Y}} + \overline{A}_{\mathcal{Y}})\left[Q_t \overline{{\mathcal{X}}}_t + \Psi_t + \varphi_t\right]+ Q_t A_\beta \beta_t + Q_t a_\mathcal{X}\right]\,dt \\
    \nonumber
    &\quad + (Q_t c_{0}-Z^\Psi_t)\,dW_t^{0} +\psi_t dt+ \dot{\varphi}_t\,dt\\
    \nonumber
    & = \Big\{ \left[\dot{{P}}_t +{P}_t A_{\mathcal{Y}} {P}_t + {P}_t A_{\mathcal{X}}\right]({\mathcal{X}}_t - \overline{{\mathcal{X}}}_t) + \left[ \dot{Q}_t + Q_t (A_{\mathcal{Y}}+ \overline{A}_{\mathcal{Y}}) Q_t + Q_t A_{\mathcal{X}} \right] \overline{{\mathcal{X}}}_t \\ 
    \nonumber
    &\quad  +\left[ \psi_t + Q_t A_\beta \beta_t + Q_t (A_{\mathcal{Y}}+ \overline{A}_{\mathcal{Y}})\Psi_t \right] + \left[\dot{\varphi}_t +Q_t a_\mathcal{X} + Q_t (A_{\mathcal{Y}}+ \overline{A}_{\mathcal{Y}})\varphi_t\right] \Big\}\,dt\\  
    \nonumber
    & \quad + \left[{P}_t C_{\mathcal{X}}({\mathcal{X}}_t - \overline{{\mathcal{X}}}_t) + {P}_t C_{\mathcal{X}} \overline{{\mathcal{X}}}_t +P_tc_K\right]\,dW_t^{K} +P_tc_X\,dW_t^{X} + (Q_t c_{0}-Z^\Psi_t)\,dW_t^{0}
\end{align}


Going back to the original expression for $\mathcal{Y}$ written in a convenient form (i.e. according to the orthogonal decomposition) and then exploiting the identities in Equation \eqref{eq:ansatz_cons}, we obtain
\begin{align*}
    d{\mathcal{Y}}_t 
    & = [B_{\mathcal{X}} ({\mathcal{X}}_t - \overline{{\mathcal{X}}}_t) + (B_{\mathcal{X}}+\overline{B}_{\mathcal{X}}) \overline{{\mathcal{X}}}_t + B_{\mathcal{Y}} ({\mathcal{Y}}_t - \overline{\mathcal{Y}}_t) + B_{\mathcal{Y}}  \overline{\mathcal{Y}}_t + B_Z ({Z}^K_t-\overline{Z}^K_t) + B_Z \overline{Z}^K_t + b_\mathcal{Y}]\,dt \\
    & \quad + {Z}^K_t\,d{W}_t^{K} + {Z}^X_t\,dW_t^{X}+ {Z}^0_t\,dW_t^{0}\\
    & = [B_{\mathcal{X}} ({\mathcal{X}}_t - \overline{{\mathcal{X}}}_t) + (B_{\mathcal{X}}+\overline{B}_{\mathcal{X}}) \overline{{\mathcal{X}}}_t + B_{\mathcal{Y}} ({\mathcal{Y}}_t - \overline{\mathcal{Y}}_t) + B_{\mathcal{Y}}  \overline{\mathcal{Y}}_t + B_Z ({Z}^K_t-\overline{Z}^K_t) + B_Z \overline{Z}^K_t + b_\mathcal{Y}]\,dt \\
    & \quad + {Z}^K_t\,d{W}_t^{K} + {Z}^X_t\,dW_t^{X}+ {Z}^0_t\,dW_t^{0}\\
    & = \Big\{B_{\mathcal{X}} ({\mathcal{X}}_t - \overline{{\mathcal{X}}}_t) + (B_{\mathcal{X}}+\overline{B}_{\mathcal{X}}) \overline{{\mathcal{X}}}_t + B_{\mathcal{Y}} {P}_t({\mathcal{X}}_t - \overline{{\mathcal{X}}}_t) + B_{\mathcal{Y}}  [Q_t \overline{{\mathcal{X}}}_t + \Psi_t + \varphi_t] \\
    & \quad + B_Z ({Z}^K_t-\overline{Z}^K_t) + B_Z \overline{Z}^K_t + b_\mathcal{Y}\Big\}\,dt + {Z}^K_t\,d{W}_t^{K} + {Z}^X_t\,dW_t^{X}+ {Z}^0_t\,dW_t^{0}.
\end{align*}
We start by matching the volatility terms and we get
\begin{align}
    {Z}^K_t = \left[{P}_t C_{\mathcal{X}}({\mathcal{X}}_t - \overline{{\mathcal{X}}}_t) + {P}_t C_{\mathcal{X}} \overline{{\mathcal{X}}}_t +P_tc_K\right],
    \quad
    {Z}^X_t = P_tc_X,
    \quad
    {Z}^0_t = (Q_t c_{0}-Z^\Psi_t).
\end{align}
Then, we can exploit the expression for  ${Z}^K_t$ and write
\begin{align}\label{eq:fin_coef_dY}
    d{\mathcal{Y}}_t 
    & = \Big\{(B_{\mathcal{X}}+B_{\mathcal{Y}} {P}_t) ({\mathcal{X}}_t - \overline{{\mathcal{X}}}_t) + (B_{\mathcal{X}}+\overline{B}_{\mathcal{X}}+   B_{\mathcal{Y}} Q_t) \overline{{\mathcal{X}}}_t +   B_{\mathcal{Y}}  \Psi_t +  B_{\mathcal{Y}}\varphi_t \\ 
    &\nonumber \quad + B_Z ({Z}^K_t-\overline{Z}^K_t) + B_Z \overline{Z}^K_t + b_\mathcal{Y}\Big\}\,dt + {Z}^K_t\,d{W}_t^{K} + {Z}^X_t\,dW_t^{X}+ {Z}^0_t\,dW_t^{0}\\  
    & \nonumber= \Big\{(B_{\mathcal{X}}+B_{\mathcal{Y}} {P}_t) ({\mathcal{X}}_t - \overline{{\mathcal{X}}}_t) + (B_{\mathcal{X}}+\overline{B}_{\mathcal{X}}+   B_{\mathcal{Y}} Q_t) \overline{{\mathcal{X}}}_t +   B_{\mathcal{Y}} \Psi_t + B_{\mathcal{Y}}\varphi_t \\  
    &\nonumber \quad + B_Z {P}_t C_{\mathcal{X}}({\mathcal{X}}_t - \overline{{\mathcal{X}}}_t) + B_Z {P}_t C_{\mathcal{X}} \overline{{\mathcal{X}}}_t + B_Z P_tc_K + b_\mathcal{Y}\Big\}\,dt + {Z}^K_t\,d{W}_t^{K} + {Z}^X_t\,dW_t^{X}+ {Z}^0_t\,dW_t^{0}\\
    &\nonumber = \Big\{(B_{\mathcal{X}}+B_{\mathcal{Y}} {P}_t  + B_Z {P}_t C_{\mathcal{X}}) ({\mathcal{X}}_t - \overline{{\mathcal{X}}}_t) + (B_{\mathcal{X}}+\overline{B}_{\mathcal{X}}+   B_{\mathcal{Y}} Q_t + B_Z {P}_t C_{\mathcal{X}} ) \overline{{\mathcal{X}}}_t + B_{\mathcal{Y}}  \Psi_t + B_{\mathcal{Y}}\varphi_t \\  
    &\nonumber \quad + B_Z P_tc_K + b_\mathcal{Y}\Big\}\,dt + {Z}^K_t\,d{W}_t^{K} + {Z}^X_t\,dW_t^{X}+ {Z}^0_t\,dW_t^{0}.
\end{align}

\begin{thebibliography}{99}
\bibitem[]{aid2023optimal} A{\"\i}d, Ren{\'e}; Biagini, Sara. Optimal dynamic regulation of carbon emissions market. Mathematical Finance 33 (2023) 80--115
\bibitem[]{carmona2013mean} Ren{\'e} Carmona; Fran{\c{c}}ois Delarue. Mean field forward-backward stochastic differential equations. Electronic Communications in Probability 18 (2013) 1 -- 15
\bibitem[]{carmona2015forward} Carmona. Forward--backward stochastic differential equations and controlled McKean--Vlasov dynamics. The Annals of Probability (2015) 2647--2700
\bibitem[]{carmona2018probabilistic2} Carmona. Probabilistic Theory of Mean Field Games with Applications II. II (2018)
\bibitem[]{del2024mean} Del Sarto, Gianmarco; Leocata, Marta; Livieri, Giulia. A Mean Field Game approach for pollution regulation of competitive firms. arXiv preprint arXiv:2407.12754 (2024)
\bibitem[]{follmer1974random} F{\"o}llmer, Hans. Random economies with many interacting agents. Journal of mathematical economics 1 (1974) 51--62
\bibitem[]{fujii2022equilibrium} Fujii, Masaaki; Takahashi, Akihiko. Equilibrium price formation with a major player and its mean field limit. ESAIM: Control, Optimisation and Calculus of Variations 28 (2022) 21
\bibitem[]{fujii2022mean} Fujii, Masaaki; Takahashi, Akihiko. A mean field game approach to equilibrium pricing with market clearing condition. SIAM Journal on Control and Optimization 60 (2022) 259--279
\bibitem[]{lacker2019mean} Lacker, Daniel; Zariphopoulou, Thaleia. Mean field and n-agent games for optimal investment under relative performance criteria. Mathematical Finance 29 (2019) 1003--1038
\bibitem[]{liptser2013statistics} Liptser, Robert S; Shiryaev, Albert N. Statistics of random processes: I. General theory. 5 (2013)
\bibitem[]{wang2018introduction} Wang, Guangchen; Wu, Zhen; Xiong, Jie; others. An introduction to optimal control of FBSDE with incomplete information. (2018)
\end{thebibliography}
\end{document}
